\problemname{Skyltar}
\noindent
Dina närmaste vänner har blivit bortskrämda från Göteborg när de fick höra om de otäcka muterade jätteråttorna, 
och har nu lyckats fly upp till det underbara Stockholm. Under sin flykt tog dina vänner hänsyn till att det märkligt nog fanns 
väldigt många vägskyltar längs vägen till Stockholm som pekade på olika små byar nere vid Lund. 
Utöver det verkade siffrorna som var givna på skyltarna lite suspiciösa, då vissa av skyltarna motsäger informationen som finns på andra skyltar.

På varje vägskylt av de alla $N$ vägskyltar som dina vänner tog hänsyn till, står det distanser till alla $M$ byar som finns i Lund. 
Siffrorna på varje vägskylt representerar avståndet till en stad avrundat till ner till närmaste heltal, då det faktiska avståndet inte nödvändigtvis var ett heltal.
Eftersom dina goda vänner är pålitliga, kan du anta att de har perfekt fotografiskt minne, vilket innebär att dina vänner minns sifforna som står på skyltarna som de såg. 
Däremot minns dina vänner inte när eller i vilken ordning som de såg dessa $N$ skyltar.

Du kan anta att Stockholm, Göteborg och de små byarna nere vid Lund ligger allting på en enda rak linje. 
Du kan dessutom anta att alla städer, byar och skyltar är så små att de kan representeras som en punkt på denna linje.
Förutom detta, kan du även anta att dina vänner har inte gått förbi någon av dessa $M$ byar.
Varken byarna eller skyltarna behöver vara på heltalspositioner. 

Givet dessa $N$ skyltar som dina vänner har sett, vad är det största mängden av dessa $N$ skyltar sådan att avståndet givna på skyltarna inte skapar några motsägelser?
Du kan tillslut även anta att åtminstone $20\%$ av alla skyltar är helt korrekta, och inte skapar någon motsägelse.

\noindent
\section*{Indata}
Första raden innehåller 


\section*{Utdata}
Skriv ut 

\section*{Poängsättning}
Din lösning kommer att testas på en mängd testfallsgrupper.
För att få poäng för en grupp så måste du klara alla testfall i gruppen.

\noindent
\begin{tabular}{| l | l | p{12cm} |}
  \hline
  \textbf{Grupp} & \textbf{Poäng} & \textbf{Gränser} \\ \hline
  $1$    & $21$         & $N \leq 15, M \leq 50$.  \\ \hline
  $2$    & $31$         & $N \leq 100, M \leq 50$ \\ \hline
  $3$    & $18$         & $N \leq 500, M \leq 50$ \\ \hline
  $4$    & $30$         & Inga ytterligare begränsningar. \\ \hline
\end{tabular}



\section*{Explanation of sample 1:}

